% Luc Destombe --- licence CC BY-NC-SA 4.0
% https://creativecommons.org/licenses/by-nc-sa/4.0/deed.fr
% Fichier autonome : compiler avec pdflatex, deux fois (signets, nombre de pages).
\documentclass[11pt,a4paper]{article}
\makeatletter
\newif\iflycee@unecolonne
\lycee@unecolonnetrue
\newif\iflycee@cm
\newif\iflycee@palatino

\RequirePackage[utf8]{inputenc}
\RequirePackage[T1]{fontenc}
\RequirePackage[french]{babel}
\RequirePackage{etoolbox}

\newcommand{\ShorthandsOff}{\@ifundefined{shorthandoff}{}{\shorthandoff{;:!?}}}
\newcommand{\ShorthandsOn}{\@ifundefined{shorthandon}{}{\shorthandon{;:!?}}}
\ShorthandsOff
\AtBeginDocument{\ShorthandsOn}
\AtBeginEnvironment{tikzpicture}{\ShorthandsOff}

\RequirePackage{geometry}
\RequirePackage{fancyhdr}
\RequirePackage{lastpage}
\RequirePackage{multicol}
\RequirePackage{array}
\RequirePackage{enumitem}
\RequirePackage{xcolor}
\RequirePackage{needspace}

\RequirePackage{amsmath,amssymb}
\iflycee@cm\else\RequirePackage{fourier}\fi
\DeclareMathAlphabet{\mathcal}{OMS}{cmsy}{m}{n}
\renewcommand{\ttdefault}{lmtt}
\begingroup\catcode`\_=\active \catcode`\^=\active
\global\let_\sb \global\let^\sp
\endgroup
\AtBeginDocument{\catcode`\_=12 \mathcode`\_="8000
  \catcode`\^=12 \mathcode`\^="8000 }
\newlength{\retraitcalculs}
\setlength{\retraitcalculs}{2em}

\RequirePackage{tikz}
\usetikzlibrary{arrows.meta,calc,babel,shapes.misc}
\RequirePackage[most]{tcolorbox}
\tcbuselibrary{skins,breakable}

\definecolor{coulDef}{HTML}{1F4E79}
\definecolor{coulProp}{HTML}{7030A0}
\definecolor{coulRem}{HTML}{C55A11}
\definecolor{coulEx}{HTML}{2E7D32}
\definecolor{coulExo}{HTML}{B03A2E}
\definecolor{coulPart}{HTML}{1F4E79}
\definecolor{coulMeth}{HTML}{1B6E4F}
\definecolor{coulCourbe}{HTML}{0B5ED7}
\definecolor{coulCourbeB}{HTML}{C00000}
\definecolor{coulCourbeC}{HTML}{2E7D32}
\definecolor{coulGrille}{HTML}{8C8C8C}
\definecolor{coulRep}{HTML}{1B6E4F}
\definecolor{coulBar}{HTML}{2E7D32}

\newcommand{\R}{\mathbb{R}}
\newcommand{\Rs}{\mathbb{R}^{*}}
\newcommand{\N}{\mathbb{N}}
\newcommand{\Ns}{\mathbb{N}^{*}}
\newcommand{\Z}{\mathbb{Z}}
\newcommand{\Q}{\mathbb{Q}}
\renewcommand{\le}{\leqslant}
\renewcommand{\ge}{\geqslant}
\renewcommand{\approx}{\simeq}
\newcommand{\vect}[1]{\overrightarrow{#1}}
\newcommand{\norme}[1]{\left\lVert #1 \right\rVert}
\newcommand{\intervalle}[2]{\left[#1\,;#2\right]}
\RequirePackage{cancel}
\renewcommand{\CancelColor}{\color{coulExo}}
\newcommand{\fois}[1]{\times{\color{coulExo}#1}}
\newcommand{\barre}[1]{\cancel{#1}}
\newlength{\lycee@eqh}\newlength{\lycee@eqd}
\newsavebox{\lycee@eqa}\newsavebox{\lycee@eqb}\newsavebox{\lycee@eqc}
\newcommand{\lycee@eqcol}[1]{\hbox{$\renewcommand{\arraystretch}{1.35}%
  \begin{array}[t]{@{}>{\displaystyle}r@{}>{\displaystyle{}}l@{}}#1\end{array}$}}
\newcommand{\lycee@eqmesure}[1]{%
  \ifdim\ht#1>\lycee@eqh\setlength{\lycee@eqh}{\ht#1}\fi
  \ifdim\dp#1>\lycee@eqd\setlength{\lycee@eqd}{\dp#1}\fi}
\newcommand{\lycee@eqtrait}{\hspace{1.2em}\tikz[baseline=0pt]\draw[dash pattern=on 2pt off 2pt,
  line width=0.5pt] (0,\lycee@eqh) -- (0,-\lycee@eqd);\hspace{1.2em}}
\newcommand{\eqproduit}[2]{\par\smallskip\noindent C'est une équation produit :
  $A\times B=0 \iff A=0$ ou $B=0$.\par\nopagebreak\smallskip
  \begingroup\setlength{\lycee@eqh}{0pt}\setlength{\lycee@eqd}{0pt}%
  \setbox\lycee@eqa=\lycee@eqcol{#1}\setbox\lycee@eqb=\lycee@eqcol{#2}%
  \lycee@eqmesure\lycee@eqa\lycee@eqmesure\lycee@eqb%
  \noindent\hspace*{\retraitcalculs}\usebox\lycee@eqa\lycee@eqtrait\usebox\lycee@eqb\endgroup\par\smallskip}
\newcommand{\eqproduittrois}[3]{\par\smallskip\noindent C'est une équation produit :
  $A\times B\times C=0 \iff A=0$ ou $B=0$ ou $C=0$.\par\nopagebreak\smallskip
  \begingroup\setlength{\lycee@eqh}{0pt}\setlength{\lycee@eqd}{0pt}%
  \setbox\lycee@eqa=\lycee@eqcol{#1}\setbox\lycee@eqb=\lycee@eqcol{#2}\setbox\lycee@eqc=\lycee@eqcol{#3}%
  \lycee@eqmesure\lycee@eqa\lycee@eqmesure\lycee@eqb\lycee@eqmesure\lycee@eqc%
  \noindent\hspace*{\retraitcalculs}\usebox\lycee@eqa\lycee@eqtrait\usebox\lycee@eqb\lycee@eqtrait\usebox\lycee@eqc\endgroup
  \par\smallskip}

\tikzset{grille/.style={coulGrille,line width=0.4pt}}
\newcommand{\quadrillage}[4]{\draw[grille,step=1] (#1,#2) grid (#3,#4);}
\tikzset{
  croix/.style={cross out,draw,line width=0.6pt,inner sep=0pt,minimum size=4pt},
  croixdroite/.style={croix,rotate=45,line width=0.8pt},
}
\newcommand{\pt}[3]{\path (#1) node[croix]{} node[#2,font=\small,inner sep=2.5pt]{$#3$};}

\newcommand{\lycee@repere}[4]{%
  \draw[grille,step=1] (#1,#2) grid (#3,#4);
  \draw[-{Stealth[length=2mm]},thick] (#1,0)--({#3+0.4},0) node[below left=-1pt]{$x$};
  \draw[-{Stealth[length=2mm]},thick] (0,#2)--(0,{#4+0.4}) node[below left=-1pt]{$y$};
  \node[below left,font=\scriptsize,inner sep=1.5pt] at (0,0) {$0$};}
\newcommand{\lycee@gradx}[1]{\foreach \k/\l in {#1}{%
  \draw (\k,0.07)--(\k,-0.07) node[below,font=\scriptsize,inner sep=1.5pt]{$\l$};}}
\newcommand{\lycee@grady}[1]{\foreach \k/\l in {#1}{%
  \draw (0.07,\k)--(-0.07,\k) node[left,font=\scriptsize,inner sep=1.5pt]{$\l$};}}
\newcommand{\lycee@intff}[2]{\left[#1\,;#2\right]}
\newcommand{\lycee@intfo}[2]{\left[#1\,;#2\right[}
\newcommand{\lycee@intof}[2]{\left]#1\,;#2\right]}
\newcommand{\lycee@intoo}[2]{\left]#1\,;#2\right[}
\newcommand{\lycee@ens}[1]{\left\{#1\right\}}
\AtBeginDocument{%
  \providecommand{\repere}{\lycee@repere}%
  \providecommand{\gradx}{\lycee@gradx}%
  \providecommand{\grady}{\lycee@grady}%
  \providecommand{\Cf}{\mathcal{C}\sb{f}}%
  \providecommand{\Cg}{\mathcal{C}\sb{g}}%
  \providecommand{\intff}{\lycee@intff}%
  \providecommand{\intfo}{\lycee@intfo}%
  \providecommand{\intof}{\lycee@intof}%
  \providecommand{\intoo}{\lycee@intoo}%
  \providecommand{\ens}{\lycee@ens}}

\newlist{questions}{enumerate}{3}
\setlist[questions,1]{label=\arabic*\textdegree),leftmargin=*,itemsep=2pt,topsep=3pt}
\setlist[questions,2]{label=\alph*),leftmargin=*,itemsep=1pt,topsep=2pt}
\setlist[questions,3]{label=\roman*),leftmargin=*,itemsep=1pt,topsep=1pt}
\setlist[itemize]{label=\textbullet,leftmargin=*,itemsep=2pt,topsep=3pt}

\newcommand{\besoinplace}[1]{\par\penalty\@M\begingroup
  \dimen@=#1\relax\dimen@ii=\pagegoal\advance\dimen@ii-\pagetotal
  \ifdim\dimen@>\dimen@ii\ifdim\dimen@ii>\z@\vfil\fi\break\fi\endgroup}

\newcommand{\lycee@pied}{\thepage/\pageref{LastPage}}
\newcommand{\entete}[2]{%
  \pagestyle{fancy}\fancyhf{}%
  \fancyhead[L]{\footnotesize\color{coulPart}\textbf{#1}}%
  \fancyhead[R]{\footnotesize\color{coulPart}\textbf{#2}}%
  \fancyfoot[C]{\footnotesize\lycee@pied}%
  \renewcommand{\headrulewidth}{0.4pt}%
  \renewcommand{\headrule}{\hbox to\headwidth{%
    \color{coulPart!40}\leaders\hrule height \headrulewidth\hfill}}}

\RequirePackage{ccicons}
\newcommand{\lycee@licence}{{\scriptsize\color{gray}%
  \copyright\ Luc Destombe --- \ccbyncsa\ CC BY-NC-SA 4.0}}
\newif\iflycee@licence \lycee@licencetrue
\newcommand{\sanslicence}{\lycee@licencefalse}
\AtBeginDocument{\iflycee@licence\AddToHookNext{shipout/foreground}{%
  \put(0,0){\rlap{\kern\dimexpr1in+\hoffset+\oddsidemargin+\textwidth\relax
    \llap{\raisebox{-\dimexpr1in+\voffset+\topmargin+\headheight+\headsep
      +\textheight+\footskip\relax}{\lycee@licence}}}}}\fi}

\newcommand{\lycee@site}{faire-des-maths.fr}
\newcommand{\lycee@version}{{\scriptsize\color{gray}\lycee@site
  \ifdefined\lyceeversion\ --- Version~\lyceeversion\fi}}
\AtBeginDocument{\iflycee@licence\AddToHookNext{shipout/foreground}{%
  \put(0,0){\rlap{\kern\dimexpr1in+\hoffset+\oddsidemargin\relax
    \raisebox{-\dimexpr1in+\voffset+\topmargin+\headheight+\headsep
      +\textheight+\footskip\relax}{\lycee@version}}}}\fi}

\newcommand{\formulesserrees}{%
  \setlength{\abovedisplayskip}{3pt}\setlength{\belowdisplayskip}{3pt}%
  \setlength{\abovedisplayshortskip}{2pt}\setlength{\belowdisplayshortskip}{3pt}}

\setlength{\parindent}{0pt}

\RequirePackage[implicit=false,hidelinks,pdfencoding=auto,psdextra,bookmarksopen,
  bookmarksopenlevel=1,pdfcreator={},pdfproducer={}]{hyperref}
\RequirePackage{bookmark}
\hypersetup{pdfauthor={Luc Destombe},pdfkeywords={CC BY-NC-SA 4.0}}
\newcounter{lycee@signet}
\newcommand{\signet}[3][]{\stepcounter{lycee@signet}%
  \edef\lycee@dest{\if\relax\detokenize{#1}\relax
    signet.\arabic{lycee@signet}\else#1\fi}%
  \hypertarget{\lycee@dest}{}\bookmark[level=#2,dest=\lycee@dest]{#3}}
\pdfstringdefDisableCommands{%
  \def\R{R}\def\N{N}\def\Z{Z}\def\Q{Q}\def\vect#1{#1}\def\Cf{Cf}\def\Cg{Cg}\def\fois#1{×#1}%
  \def\barre#1{#1}\def\intervalle#1#2{[#1;#2]}\def\intff#1#2{[#1;#2]}%
  \def\textsuperscript#1{#1}\def\dfrac#1#2{#1/#2}\def\frac#1#2{#1/#2}%
  \def\vec#1{#1⃗}}%
\RequirePackage[official]{eurosym}

\tcbset{exercice corrige/.style={
  enhanced,breakable,before upper=\raggedright,
  colback=white,colframe=coulExo,
  boxrule=0.8pt,arc=2pt,
  left=6pt,right=6pt,top=8pt,bottom=5pt,
  before skip=9pt,after skip=4pt,
  coltitle=white,fonttitle=\bfseries\small,
  attach boxed title to top left={xshift=8pt,yshift=-\tcboxedtitleheight/2},
  boxed title style={colback=coulExo,arc=2pt,boxrule=0pt}}}
\newcounter{exo}
\newtcolorbox{exercice}[1][]{exercice corrige,
  phantom={\signet[exercice-\arabic{exo}]{2}{Exercice \arabic{exo}}},
  title={Exercice \theexo\ifx\relax#1\relax\else{} --- #1\fi}}
\newenvironment{exo}[1][\relax]{\refstepcounter{exo}\begin{exercice}[#1]}{\end{exercice}}

\geometry{top=1.8cm,bottom=1cm,left=1cm,right=1cm,headheight=14pt,footskip=0.6cm}
\renewcommand{\lycee@pied}{\thepage}

\renewcommand{\theexo}{\ifnum\value{exo}<10 0\fi\arabic{exo}}
\tcbset{exercice corrige/.style={
  enhanced,breakable,before upper=\raggedright,
  colback=white,colframe=coulExo,
  boxrule=0.9pt,arc=2pt,
  left=8pt,right=8pt,top=8pt,bottom=6pt,
  before skip=8pt,after skip=6pt,
  coltitle=white,fonttitle=\bfseries,
  attach boxed title to top left={xshift=8pt,yshift=-\tcboxedtitleheight/2},
  boxed title style={colback=coulExo,arc=2pt,boxrule=0pt}}}

\newcommand{\entetefiche}[2]{%
  \begin{center}
    \begin{tcolorbox}[enhanced,width=0.92\linewidth,colback=white,colframe=coulPart,
        boxrule=1.4pt,arc=3pt,halign=center,top=6pt,bottom=6pt]
      {\large\bfseries\color{coulPart} CORRIGÉ --- FICHE D'EXERCICES}\\[3pt]
      {\normalsize\bfseries #1}\\[2pt]
      {\normalsize #2}
    \end{tcolorbox}
  \end{center}}

\newcounter{partie}
\newcommand{\partie}[1]{%
  \stepcounter{partie}%
  \besoinplace{8\baselineskip}\vspace{4pt}%
  \begin{tcolorbox}[enhanced,colback=coulPart!8,colframe=coulPart,boxrule=0pt,
      leftrule=3pt,arc=0pt,left=6pt,right=4pt,top=3pt,bottom=3pt,
      before skip=6pt,after skip=2pt]
    \bfseries\color{coulPart}\Roman{partie} --- #1\signet{1}{\Roman{partie} --- #1}
  \end{tcolorbox}}
\newcommand{\partiesansnum}[1]{%
  \besoinplace{8\baselineskip}\vspace{4pt}%
  \begin{tcolorbox}[enhanced,colback=coulPart!8,colframe=coulPart,boxrule=0pt,
      leftrule=3pt,arc=0pt,left=6pt,right=4pt,top=3pt,bottom=3pt,
      before skip=6pt,after skip=2pt]
    \bfseries\color{coulPart}#1\signet{1}{#1}
  \end{tcolorbox}}

\setlist[questions,1]{label=\arabic*\textdegree),leftmargin=*,itemsep=3pt,topsep=3pt}
\setlist[questions,2]{label=\alph*),leftmargin=*,itemsep=2pt,topsep=2pt}
\setlist[questions,3]{label=\roman*),leftmargin=*,itemsep=1pt,topsep=1pt}
\newlist{expr}{enumerate}{1}
\setlist[expr]{label={\Alph*\ $=$},leftmargin=*,itemsep=2pt,topsep=3pt}

\newcommand{\res}[1]{{\color{coulRep}#1}}
\newcommand{\rem}[1]{\par\smallskip{\small\itshape\color{black!60}$\rightarrow$ #1}}
\newenvironment{calculs}{\par\smallskip\noindent$\displaystyle\begin{aligned}[t]}%
  {\end{aligned}$\par\smallskip}
\newcommand{\justif}[1]{\quad\text{\small\itshape\color{black!60}#1}}
\newcommand{\cas}[1]{\par\smallskip\noindent\textbf{\color{coulExo}#1}\vadjust{\nobreak}}
\newcommand{\explique}[2]{\par\smallskip\noindent
  \begin{minipage}[t]{0.6\linewidth}\raggedright #1\end{minipage}\hfill
  \begin{minipage}[t]{0.37\linewidth}\raggedright\leavevmode
    {\small\itshape\color{black!60}#2}\end{minipage}\par}

\setlength{\parskip}{2pt}

\makeatother
\entete{Chapitre 7 --- Variables aléatoires, loi de Bernoulli --- corrigé}{1\textsuperscript{re} STMG}
\let\casligne\cas
\renewcommand{\cas}[1]{\casligne{#1}\par\nopagebreak}
\let\partieligne\partie
\renewcommand{\partie}[1]{\besoinplace{16\baselineskip}\partieligne{#1}}
\BeforeBeginEnvironment{exo}{\besoinplace{10\baselineskip}}
\setlength{\parindent}{0pt}
\setlength{\jot}{5pt}

\tikzset{evt/.style={inner sep=2pt},
  proba/.style={font=\footnotesize,fill=white,inner sep=1pt}}
\newcommand{\feuille}[2]{\node[right,font=\footnotesize] at (#1.east) {#2};}
\newcommand{\arbredeux}[3]{%
  \node (O) at (0,0) {};
  \node[evt] (S) at (2.2,1.2) {$#1$};
  \node[evt] (E) at (2.2,-1.2) {$\overline{#1}$};
  \node[evt] (SS) at (4.4,1.8) {$#1$};
  \node[evt] (SE) at (4.4,0.6) {$\overline{#1}$};
  \node[evt] (ES) at (4.4,-0.6) {$#1$};
  \node[evt] (EE) at (4.4,-1.8) {$\overline{#1}$};
  \draw (O.center) -- node[proba]{$#2$} (S);
  \draw (O.center) -- node[proba]{$#3$} (E);
  \draw (S) -- node[proba]{$#2$} (SS);
  \draw (S) -- node[proba]{$#3$} (SE);
  \draw (E) -- node[proba]{$#2$} (ES);
  \draw (E) -- node[proba]{$#3$} (EE);}
\newcommand{\arbretrois}[3]{%
  \node (O) at (0,0) {};
  \node[evt] (S) at (1.9,2) {$#1$};
  \node[evt] (E) at (1.9,-2) {$\overline{#1}$};
  \foreach \a/\y in {S/2,E/-2}{%
    \node[evt] (\a S) at (3.8,\y+1) {$#1$};
    \node[evt] (\a E) at (3.8,\y-1) {$\overline{#1}$};
    \draw (\a) -- node[proba]{$#2$} (\a S);
    \draw (\a) -- node[proba]{$#3$} (\a E);}
  \foreach \a/\y in {SS/3,SE/1,ES/-1,EE/-3}{%
    \node[evt] (\a S) at (5.7,\y+0.5) {$#1$};
    \node[evt] (\a E) at (5.7,\y-0.5) {$\overline{#1}$};
    \draw (\a) -- node[proba,pos=0.62]{$#2$} (\a S);
    \draw (\a) -- node[proba,pos=0.62]{$#3$} (\a E);}
  \draw (O.center) -- node[proba]{$#2$} (S);
  \draw (O.center) -- node[proba]{$#3$} (E);}
\newcommand{\arbre}[8]{%
  \node (O) at (0,0) {};
  \node[evt] (A) at (2.6,1.4) {$#1$};
  \node[evt] (nA) at (2.6,-1.4) {$\overline{#1}$};
  \node[evt] (AB) at (5.2,2.1) {$#2$};
  \node[evt] (AnB) at (5.2,0.7) {$\overline{#2}$};
  \node[evt] (nAB) at (5.2,-0.7) {$#2$};
  \node[evt] (nAnB) at (5.2,-2.1) {$\overline{#2}$};
  \draw (O.center) -- node[proba]{$#3$} (A);
  \draw (O.center) -- node[proba]{$#4$} (nA);
  \draw (A) -- node[proba]{$#5$} (AB);
  \draw (A) -- node[proba]{$#6$} (AnB);
  \draw (nA) -- node[proba]{$#7$} (nAB);
  \draw (nA) -- node[proba]{$#8$} (nAnB);}
\newtcblisting{python}{%
  listing only,enhanced,breakable,
  colback=black!3,colframe=black!35,boxrule=0.5pt,arc=1pt,
  left=6pt,right=6pt,top=2pt,bottom=2pt,
  listing options={language=Python,basicstyle=\ttfamily\small,
    keywordstyle=\color{coulMeth}\bfseries,columns=fullflexible,
    keepspaces=true,showstringspaces=false,aboveskip=0pt,belowskip=0pt}}

\begin{document}

\entetefiche{Chapitre 7 --- Variables aléatoires, loi de Bernoulli}{Première STMG}

\partie{Variable aléatoire}

\begin{exo}[Calculer des probabilités avec une variable aléatoire]
\cas{1)}
\begin{center}
\renewcommand{\arraystretch}{1.4}
\begin{tabular}{|l|c|c|c|c|c|c|}
\hline
Issue & $1$ & $2$ & $3$ & $4$ & $5$ & $6$\\
\hline
Valeur de $X$ & $-2$ & $-2$ & $-2$ & $0$ & $3$ & $3$\\
\hline
\end{tabular}
\end{center}
\explique{$X$ prend les valeurs \res{$-2$ ; $0$ et $3$}.}{perte de $2$~€, rien, ou gain de
$3$~€}
\cas{2)}
\explique{Le dé est équilibré : les $6$ issues ont la même probabilité.}{équiprobabilité}
\begin{calculs}
P(X=3) &= \dfrac{2}{6} &&\justif{issues $5$ et $6$, soit $2$ issues sur $6$}\\
P(X=3) &= \dfrac{\barre{2}\times 1}{\barre{2}\times 3} &&\justif{on simplifie par $2$}\\
P(X=3) &= \res{\dfrac{1}{3}} &&
\end{calculs}
\begin{calculs}
P(X=-2) &= \dfrac{3}{6} &&\justif{issues $1$ ; $2$ ; $3$}\\
P(X=-2) &= \dfrac{\barre{3}\times 1}{\barre{3}\times 2} &&\justif{on simplifie par $3$}\\
P(X=-2) &= \res{\dfrac{1}{2}} &&
\end{calculs}
\cas{3)}
\explique{$\{X\ge 0\}$ : \res{\og on ne perd pas d'argent\fg{}}, c'est-à-dire \og on obtient
un $4$, un $5$ ou un $6$\fg{}.}{$X$ vaut $0$ ou $3$}
\begin{calculs}
P(X\ge 0) &= P(X=0)+P(X=3) &&\justif{les valeurs $0$ et $3$}\\
P(X\ge 0) &= \dfrac{1}{6}+\dfrac{2}{6} &&\justif{issue $4$, puis issues $5$ et $6$}\\
P(X\ge 0) &= \dfrac{3}{6} &&\\
P(X\ge 0) &= \res{\dfrac{1}{2}} &&\justif{on simplifie par $3$}
\end{calculs}
\cas{4)}
\begin{calculs}
P(X\le 0) &= 1-P(X=3) &&\justif{contraire de $\{X=3\}$ : on n'obtient ni $5$ ni $6$}\\
P(X\le 0) &= 1-\dfrac{1}{3} &&\justif{résultat de la question 2}\\
P(X\le 0) &= \res{\dfrac{2}{3}} &&
\end{calculs}
\rem{On peut aussi compter les issues $1$ ; $2$ ; $3$ ; $4$, soit $\frac{4}{6}$, qui se simplifie
en $\frac{2}{3}$.}
\end{exo}

\begin{exo}[En phrase complète]
\cas{1)}
\explique{a) \res{$\{X=7\}$}}{exactement : une seule valeur}
\explique{b) \res{$\{X\le 3\}$}}{au plus $3$, c'est $3$ ou moins}
\explique{c) \res{$\{X\ge 8\}$}}{au moins $8$, c'est $8$ ou plus}
\explique{d) \res{$\{X<2\}$}}{moins de $2$, c'est $0$ ou $1$, sans $2$}
\explique{e) \res{$\{X>6\}$}}{plus de $6$, c'est $7$ ou plus, sans $6$}
\explique{f) \res{$\{X=10\}$}}{les $10$ clients rencontrés}
\rem{\og Au plus\fg{} et \og au moins\fg{} comprennent le nombre ; \og moins de\fg{} et
\og plus de\fg{} l'excluent.}
\cas{2)}
\explique{$\{X=0\}$ : \og \res{aucun client ne passe de commande}\fg{}.}{$X$ vaut $0$}
\explique{$\{X\le 9\}$ : \og \res{au plus $9$ clients passent une commande}\fg{},
c'est-à-dire \og au moins un client ne passe pas de commande\fg{}.}{toutes les valeurs sauf $10$}
\explique{$\{X\ge 1\}$ : \og \res{au moins un client passe une commande}\fg{}.}{toutes les valeurs
sauf $0$}
\cas{3)}
\explique{\og Entre $4$ et $6$, bornes comprises\fg{} : $X$ vaut $4$, $5$ ou $6$, donc la
probabilité vaut \res{$P(X=4)+P(X=5)+P(X=6)$}.}%
{on ajoute les probabilités des valeurs qui réalisent l'événement}
\end{exo}

\begin{exo}[Des journées promotion]
\cas{1)}%
\begin{center}
\renewcommand{\arraystretch}{1.4}
\begin{tabular}{|l|c|c|c|c|c|c|c|c|}
\hline
Issue & $PPP$ & $PPN$ & $PNP$ & $PNN$ & $NPP$ & $NPN$ & $NNP$ & $NNN$\\
\hline
Valeur de $X$ & $3$ & $2$ & $2$ & $1$ & $2$ & $1$ & $1$ & $0$\\
\hline
\end{tabular}
\end{center}
\explique{On écrit d'abord les issues qui commencent par $P$, puis celles qui commencent par
$N$ : rien n'est oublié.}{$2$ choix pour chaque jour : $2\times 2\times 2=8$ issues}
\cas{2)}
\explique{$X$ prend les valeurs \res{$0$ ; $1$ ; $2$ et $3$}.}{d'aucune journée promotion à
trois}
\cas{3)}
\begin{calculs}
P(X=2) &= \res{\dfrac{3}{8}} &&\justif{issues $PPN$ ; $PNP$ ; $NPP$, soit $3$ issues sur
$8$}
\end{calculs}
\begin{calculs}
P(X=0) &= \res{\dfrac{1}{8}} &&\justif{la seule issue $NNN$}
\end{calculs}
\cas{4)}
\explique{$\{X\ge 1\}$ : \og \res{il y a au moins une journée promotion}\fg{}.}{son contraire
est $\{X=0\}$, l'issue $NNN$}
\begin{calculs}
P(X\ge 1) &= 1-P(X=0) &&\justif{événement contraire}\\
P(X\ge 1) &= 1-\dfrac{1}{8} &&\\
P(X\ge 1) &= \res{\dfrac{7}{8}} &&
\end{calculs}
\end{exo}

\begin{exo}[À l'ouverture d'une boutique]
\cas{1)}
\explique{Multiple de $5$ ($5$ et $10$) : gain $5-1=4$ ; multiple de $3$ ($3$ ; $6$ ; $9$) :
gain $1-1=0$ ; sinon : gain $-1$.}{on déduit le prix de la partie}
\begin{center}
\renewcommand{\arraystretch}{1.4}
\begin{tabular}{|l|c|c|c|c|c|c|c|c|c|c|}
\hline
Numéro & $1$ & $2$ & $3$ & $4$ & $5$ & $6$ & $7$ & $8$ & $9$ & $10$\\
\hline
Valeur de $X$ & $-1$ & $-1$ & $0$ & $-1$ & $4$ & $0$ & $-1$ & $-1$ & $0$ & $4$\\
\hline
\end{tabular}
\end{center}
\explique{$X$ prend les valeurs \res{$-1$ ; $0$ et $4$}.}{aucun numéro n'est à la fois
multiple de $5$ et de $3$ (il faudrait $15$)}
\cas{2)}
\explique{Les $10$ boules ont la même probabilité d'être tirées.}{tirage au hasard}
\begin{calculs}
P(X=4) &= \dfrac{2}{10} &&\justif{numéros $5$ et $10$}\\
P(X=4) &= \res{\dfrac{1}{5}} &&\justif{on simplifie par $2$ (soit $0{,}2$)}
\end{calculs}
\begin{calculs}
P(X=-1) &= \dfrac{5}{10} &&\justif{numéros $1$ ; $2$ ; $4$ ; $7$ ; $8$}\\
P(X=-1) &= \res{\dfrac{1}{2}} &&\justif{on simplifie par $5$}
\end{calculs}
\cas{3)}
\explique{$\{X\ge 0\}$ : \og \res{on ne perd pas d'argent}\fg{}, c'est-à-dire \og le numéro est
un multiple de $3$ ou de $5$\fg{}.}{$X$ vaut $0$ ou $4$}
\begin{calculs}
P(X\ge 0) &= P(X=0)+P(X=4) &&\justif{les valeurs $0$ et $4$}\\
P(X\ge 0) &= \dfrac{3}{10}+\dfrac{2}{10} &&\justif{numéros $3$ ; $6$ ; $9$, puis $5$ ; $10$}\\
P(X\ge 0) &= \dfrac{5}{10} &&\\
P(X\ge 0) &= \res{\dfrac{1}{2}} &&
\end{calculs}
\cas{4)}
\begin{calculs}
P(X<4) &= 1-P(X=4) &&\justif{contraire de $\{X<4\}$ : $\{X\ge 4\}$, soit $\{X=4\}$}\\
P(X<4) &= 1-\dfrac{1}{5} &&\justif{question 2}\\
P(X<4) &= \res{\dfrac{4}{5}} &&\justif{soit $0{,}8$}
\end{calculs}
\end{exo}

\partie{Loi de probabilité d'une variable aléatoire}

\begin{exo}[Déterminer et utiliser une loi]
\cas{a)}
\explique{Avec un lot de $50$~€, le gain vaut $50-3=47$ ; avec un lot de $10$~€, $10-3=7$ ;
sans lot, $-3$. Donc $X$ prend les valeurs $-3$ ; $7$ et $47$.}{on déduit le prix du
ticket}
\begin{calculs}
P(X=47) &= \dfrac{2}{200} = \res{0{,}01} &&\justif{$2$ tickets sur $200$}\\
P(X=7) &= \dfrac{10}{200} = \res{0{,}05} &&\justif{$10$ tickets sur $200$}\\
P(X=-3) &= \dfrac{188}{200} = \res{0{,}94} &&\justif{$200-2-10=188$ tickets sans lot}
\end{calculs}
\begin{center}
\renewcommand{\arraystretch}{1.4}
\begin{tabular}{|c|c|c|c|}
\hline
$x_i$ & $-3$ & $7$ & $47$\\
\hline
$P(X=x_i)$ & $0{,}94$ & $0{,}05$ & $0{,}01$\\
\hline
\end{tabular}
\end{center}
\explique{Vérification : $0{,}94+0{,}05+0{,}01=1$.}{propriété 1}
\cas{b)}%
\begin{minipage}[c]{0.58\linewidth}\raggedright
Le tableau donne le plus petit des deux nombres pour chacune des $4\times 4=16$ issues, qui
ont la même probabilité $\frac{1}{16}$. $Y$ prend les valeurs $1$ ; $2$ ; $3$ et $4$ ; on
compte les cases de chaque valeur (par exemple, $Y=1$ dans $7$ cases : la première ligne et
la première colonne).
\end{minipage}\hfill
\begin{minipage}[c]{0.38\linewidth}\centering
\renewcommand{\arraystretch}{1.3}
\begin{tabular}{|c|c|c|c|c|}
\hline
 & $1$ & $2$ & $3$ & $4$\\
\hline
$1$ & $1$ & $1$ & $1$ & $1$\\
\hline
$2$ & $1$ & $2$ & $2$ & $2$\\
\hline
$3$ & $1$ & $2$ & $3$ & $3$\\
\hline
$4$ & $1$ & $2$ & $3$ & $4$\\
\hline
\end{tabular}
\end{minipage}
\begin{center}
\renewcommand{\arraystretch}{2.4}
\begin{tabular}{|c|c|c|c|c|}
\hline
$y_i$ & $1$ & $2$ & $3$ & $4$\\
\hline
$P(Y=y_i)$ & $\dfrac{7}{16}$ & $\dfrac{5}{16}$ & $\dfrac{3}{16}$ & $\dfrac{1}{16}$\\
\hline
\end{tabular}
\end{center}
\explique{Vérification : $7+5+3+1=16$, donc la somme des probabilités vaut
$\frac{16}{16}=1$.}{on garde le dénominateur $16$ pour vérifier}
\cas{c)}
\begin{calculs}
0{,}35+0{,}30+a+0{,}10+0{,}05 &= 1 &&\justif{la somme des probabilités vaut $1$}\\
0{,}80+a &= 1 &&\\
a &= 1-0{,}80 &&\\
a &= \res{0{,}20} &&
\end{calculs}
\begin{calculs}
P(Z\le 1) &= P(Z=0)+P(Z=1) &&\justif{les valeurs $0$ et $1$}\\
P(Z\le 1) &= 0{,}35+0{,}30 &&\\
P(Z\le 1) &= \res{0{,}65} &&
\end{calculs}
\explique{\og Au moins $2$ clients attendent\fg{} s'écrit $\{Z\ge 2\}$.}{$Z$ vaut $2$ ; $3$
ou $4$ : c'est le contraire de $\{Z\le 1\}$}
\begin{calculs}
P(Z\ge 2) &= 1-P(Z\le 1) &&\justif{événement contraire}\\
P(Z\ge 2) &= 1-0{,}65 &&\\
P(Z\ge 2) &= \res{0{,}35} &&\justif{on vérifie : $0{,}20+0{,}10+0{,}05=0{,}35$}
\end{calculs}
\end{exo}

\begin{exo}[Des réclamations]
\cas{1)}
\begin{calculs}
0{,}05+0{,}15+0{,}30+a+0{,}15+0{,}10 &= 1 &&\justif{la somme des probabilités vaut $1$}\\
0{,}75+a &= 1 &&\\
a &= 1-0{,}75 &&\\
a &= \res{0{,}25} &&
\end{calculs}
\cas{2)}
\explique{$\{X\le 2\}$ : \og \res{le service reçoit au plus $2$ réclamations}\fg{}.}{$X$ vaut $0$ ;
$1$ ou $2$}
\begin{calculs}
P(X\le 2) &= 0{,}05+0{,}15+0{,}30 &&\justif{les valeurs $0$ ; $1$ ; $2$}\\
P(X\le 2) &= \res{0{,}50} &&
\end{calculs}
\explique{$\{X\ge 4\}$ : \og \res{le service reçoit au moins $4$ réclamations}\fg{}.}{$X$ vaut $4$
ou $5$}
\begin{calculs}
P(X\ge 4) &= 0{,}15+0{,}10 &&\justif{les valeurs $4$ et $5$}\\
P(X\ge 4) &= \res{0{,}25} &&
\end{calculs}
\cas{3)}
\explique{\og Plus de $2$ réclamations\fg{} s'écrit $\{X>2\}$, c'est-à-dire $\{X\ge 3\}$ : c'est
le contraire de $\{X\le 2\}$.}{$2$ exclu}
\begin{calculs}
P(X>2) &= 1-P(X\le 2) &&\justif{événement contraire}\\
P(X>2) &= 1-0{,}50 &&\justif{question 2}\\
P(X>2) &= \res{0{,}50} &&\justif{on vérifie : $0{,}25+0{,}15+0{,}10=0{,}50$}
\end{calculs}
\cas{4)}
\begin{calculs}
0{,}2+0{,}3+0{,}3+0{,}3 &= 1{,}1 &&\justif{somme des probabilités du tableau}
\end{calculs}
\explique{La somme vaut $1{,}1$ et non $1$ : \res{ce tableau ne peut pas être une loi de
probabilité}.}{propriété 1 : la somme des probabilités d'une loi vaut toujours $1$}
\end{exo}

\begin{exo}[Un dé publicitaire]
\cas{1)}
\explique{Chaque face paire a deux fois plus de chances de sortir qu'une face impaire :
la probabilité d'obtenir $2$ vaut \res{$2p$}.}{de même pour $4$ et pour $6$}
\cas{2)}
\begin{calculs}
p+2p+p+2p+p+2p &= 1 &&\justif{faces $1$ à $6$ : la somme des probabilités vaut $1$}\\
9p &= 1 &&\justif{$3$ fois $p$ et $3$ fois $2p$}\\
p &= \res{\dfrac{1}{9}} &&\justif{on divise par $9$}
\end{calculs}
\cas{3)}%
\explique{Faces impaires : $p=\frac{1}{9}$ ; faces paires : $2p=\frac{2}{9}$.}{question 1}
\begin{center}
\renewcommand{\arraystretch}{2.4}
\begin{tabular}{|c|c|c|c|c|c|c|}
\hline
$x_i$ & $1$ & $2$ & $3$ & $4$ & $5$ & $6$\\
\hline
$P(X=x_i)$ & $\dfrac{1}{9}$ & $\dfrac{2}{9}$ & $\dfrac{1}{9}$ & $\dfrac{2}{9}$
  & $\dfrac{1}{9}$ & $\dfrac{2}{9}$\\
\hline
\end{tabular}
\end{center}
\explique{Vérification : $1+2+1+2+1+2=9$, donc la somme vaut $\frac{9}{9}=1$.}{propriété 1}
\cas{4)}
\begin{calculs}
P(X\ge 4) &= P(X=4)+P(X=5)+P(X=6) &&\justif{les valeurs $4$ ; $5$ ; $6$}\\
P(X\ge 4) &= \dfrac{2}{9}+\dfrac{1}{9}+\dfrac{2}{9} &&\\
P(X\ge 4) &= \res{\dfrac{5}{9}} &&
\end{calculs}
\begin{calculs}
P(\text{pair}) &= P(X=2)+P(X=4)+P(X=6) &&\justif{les faces $2$ ; $4$ ; $6$}\\
P(\text{pair}) &= \dfrac{2}{9}+\dfrac{2}{9}+\dfrac{2}{9} &&\\
P(\text{pair}) &= \dfrac{6}{9} &&\\
P(\text{pair}) &= \dfrac{\barre{3}\times 2}{\barre{3}\times 3} &&\justif{on simplifie par $3$}\\
P(\text{pair}) &= \res{\dfrac{2}{3}} &&
\end{calculs}
\rem{Le dé n'est pas équilibré : on ne peut pas compter les faces favorables et diviser
par $6$ ; on additionne les probabilités de la loi.}
\end{exo}

\partie{Espérance}

\begin{exo}[Calculer et interpréter une espérance]
\cas{a)}
\begin{calculs}
E(X) &= 0\times 0{,}10+1\times 0{,}30+2\times 0{,}40+3\times 0{,}20 &&\justif{valeur
$\times$ probabilité}\\
E(X) &= 0+0{,}30+0{,}80+0{,}60 &&\\
E(X) &= \res{1{,}7} &&
\end{calculs}
\explique{Sur un grand nombre d'heures, le drone \res{livre en moyenne $1{,}7$ colis par
heure}.}{propriété 2 : une moyenne n'est pas forcément entière}
\cas{b)}
\begin{calculs}
E(Y) &= (-2)\times 0{,}85+3\times 0{,}10+18\times 0{,}05 &&\justif{valeur $\times$
probabilité}\\
E(Y) &= -1{,}70+0{,}30+0{,}90 &&\\
E(Y) &= \res{-0{,}5} &&
\end{calculs}
\explique{Sur un grand nombre de tickets, un visiteur \res{perd en moyenne $0{,}50$~€ par
ticket} ; l'association \res{gagne en moyenne $0{,}50$~€ par ticket vendu}.}{$E(Y)<0$ :
perte pour le joueur, gain pour l'organisateur}
\end{exo}

\begin{exo}[Deux stands]
\cas{1)}
\begin{calculs}
E(X) &= (-1)\times 0{,}75+2\times 0{,}20+5\times 0{,}05 &&\justif{valeur $\times$
probabilité}\\
E(X) &= -0{,}75+0{,}40+0{,}25 &&\\
E(X) &= \res{-0{,}10} &&
\end{calculs}
\begin{calculs}
E(Y) &= (-2)\times 0{,}85+4\times 0{,}10+20\times 0{,}05 &&\justif{valeur $\times$
probabilité}\\
E(Y) &= -1{,}70+0{,}40+1{,}00 &&\\
E(Y) &= \res{-0{,}30} &&
\end{calculs}
\cas{2)}
\explique{Sur un grand nombre de parties, on perd en moyenne $0{,}10$~€ par partie au stand A
et $0{,}30$~€ au stand B : \res{le stand A est le moins défavorable}.}{$-0{,}10>-0{,}30$}
\cas{3)}
\begin{calculs}
P(X>0) &= P(X=2)+P(X=5) &&\justif{les gains positifs}\\
P(X>0) &= 0{,}20+0{,}05 &&\\
P(X>0) &= \res{0{,}25} &&
\end{calculs}
\begin{calculs}
P(Y>0) &= P(Y=4)+P(Y=20) &&\justif{les gains positifs}\\
P(Y>0) &= 0{,}10+0{,}05 &&\\
P(Y>0) &= \res{0{,}15} &&
\end{calculs}
\cas{4)}
\explique{\res{L'ami a tort} : au stand B, on ne gagne $20$~€ qu'avec la probabilité
$0{,}05$, on gagne moins souvent qu'au stand A ($0{,}15$ contre $0{,}25$) et on perd en
moyenne davantage ($0{,}30$~€ contre $0{,}10$~€ par partie).}{un gros lot rare ne fait pas
un bon jeu : on regarde l'espérance}
\end{exo}

\begin{exo}[Au magasin de cycles]
\cas{1)}%
\begin{minipage}[c]{0.50\linewidth}\raggedright
$P(V)=0{,}2$, donc $P(\overline{V})=0{,}8$ ; $P_V(C)=0{,}6$, donc
$P_V(\overline{C})=0{,}4$ ; $P_{\overline{V}}(C)=0{,}1$, donc
$P_{\overline{V}}(\overline{C})=0{,}9$.

\smallskip
Au bout de chaque chemin : la somme dépensée.
\end{minipage}\hfill
\begin{minipage}[c]{0.46\linewidth}\centering
\begin{tikzpicture}
  \arbre{V}{C}{0{,}2}{0{,}8}{0{,}6}{0{,}4}{0{,}1}{0{,}9}
  \feuille{AB}{$X=1\,250$}
  \feuille{AnB}{$X=1\,200$}
  \feuille{nAB}{$X=50$}
  \feuille{nAnB}{$X=0$}
\end{tikzpicture}
\end{minipage}
\cas{2)}
\explique{$X$ prend les valeurs \res{$0$ ; $50$ ; $1\,200$ et $1\,250$}.}{rien ; un casque
seul ; un vélo seul ; un vélo et un casque ($1\,200+50$)}
\cas{3)}
\begin{calculs}
P(X=1\,250) &= 0{,}2\times 0{,}6 = \res{0{,}12} &&\justif{chemin $V$ puis $C$}\\
P(X=1\,200) &= 0{,}2\times 0{,}4 = \res{0{,}08} &&\justif{chemin $V$ puis $\overline{C}$}\\
P(X=50) &= 0{,}8\times 0{,}1 = \res{0{,}08} &&\justif{chemin $\overline{V}$ puis $C$}\\
P(X=0) &= 0{,}8\times 0{,}9 = \res{0{,}72} &&\justif{chemin $\overline{V}$ puis
$\overline{C}$}
\end{calculs}
\begin{center}
\renewcommand{\arraystretch}{1.4}
\begin{tabular}{|c|c|c|c|c|}
\hline
$x_i$ & $0$ & $50$ & $1\,200$ & $1\,250$\\
\hline
$P(X=x_i)$ & $0{,}72$ & $0{,}08$ & $0{,}08$ & $0{,}12$\\
\hline
\end{tabular}
\end{center}
\explique{Vérification : $0{,}72+0{,}08+0{,}08+0{,}12=1$.}{propriété 1}
\cas{4)}
\begin{calculs}
E(X) &= 0\times 0{,}72+50\times 0{,}08+1\,200\times 0{,}08+1\,250\times 0{,}12
  &&\justif{valeur $\times$ probabilité}\\
E(X) &= 0+4+96+150 &&\\
E(X) &= \res{250} &&
\end{calculs}
\explique{Sur un grand nombre de clients, un client \res{dépense en moyenne $250$~€} en vélos
et en casques.}{la plupart ne dépensent rien, quelques-uns plus de $1\,000$~€}
\cas{5)}
\begin{calculs}
\text{chiffre d'affaires} &= 400\times 250 &&\justif{nombre de clients $\times$ dépense
moyenne}\\
\text{chiffre d'affaires} &= \res{100\,000} &&
\end{calculs}
\explique{Le magasin peut prévoir \res{environ $100\,000$~€ par mois} pour les vélos et les
casques.}{une prévision : le chiffre effectif fluctue d'un mois à l'autre}
\end{exo}

\begin{exo}[Le coût d'une trottinette]
\cas{1)}
\explique{$X$ prend les valeurs \res{$200$ ; $220$ ; $230$ et $250$}.}{$200$ sans défaut,
$200+20$, $200+30$, $200+50$}
\cas{2)}%
\begin{center}
\renewcommand{\arraystretch}{1.4}
\begin{tabular}{|c|c|c|c|c|}
\hline
$x_i$ & $200$ & $220$ & $230$ & $250$\\
\hline
$P(X=x_i)$ & $0{,}92$ & $0{,}04$ & $0{,}03$ & $0{,}01$\\
\hline
\end{tabular}
\end{center}
\explique{Vérification : $0{,}92+0{,}04+0{,}03+0{,}01=1$.}{les pourcentages de l'énoncé,
écrits en décimal}
\cas{3)}
\begin{calculs}
E(X) &= 200\times 0{,}92+220\times 0{,}04+230\times 0{,}03+250\times 0{,}01
  &&\justif{valeur $\times$ probabilité}\\
E(X) &= 184+8{,}8+6{,}9+2{,}5 &&\\
E(X) &= \res{202{,}2} &&
\end{calculs}
\explique{Sur un grand nombre de trottinettes, une trottinette \res{coûte en moyenne
$202{,}20$~€} à l'entreprise, réparation comprise.}{propriété 2}
\cas{4)}
\explique{Elle dépense en moyenne $202{,}20$~€ et reçoit $200$~€ : \res{elle perd en moyenne
$2{,}20$~€ par trottinette}, donc elle ne gagne pas d'argent.}{$200-202{,}2=-2{,}2$}
\cas{5)}
\begin{calculs}
\text{prix} &= 202{,}2+50 &&\justif{coût moyen + gain moyen voulu}\\
\text{prix} &= \res{252{,}2} &&
\end{calculs}
\explique{Elle doit vendre chaque trottinette \res{$252{,}20$~€}.}{elle gagne alors en moyenne
$252{,}2-202{,}2=50$~€}
\end{exo}

\partie{Loi de Bernoulli}

\begin{exo}[Reconnaître une loi de Bernoulli]
\cas{a)}
\explique{Le passage a deux issues ; le succès est \og le bus arrive à l'heure\fg{}, de
probabilité $0{,}85$. Donc \res{$X$ suit la loi de Bernoulli de paramètre
$p=0{,}85$}.}{$X$ vaut $1$ (succès) ou $0$ (échec)}
\begin{center}
\renewcommand{\arraystretch}{1.4}
\begin{tabular}{|c|c|c|}
\hline
$x_i$ & $1$ & $0$\\
\hline
$P(X=x_i)$ & $0{,}85$ & $0{,}15$\\
\hline
\end{tabular}
\end{center}
\explique{$E(X)=p$, soit \res{$E(X)=0{,}85$}.}{propriété 3}
\cas{b)}
\explique{Le succès est \og obtenir un multiple de $3$\fg{}, c'est-à-dire $3$ ; $6$ ; $9$ ou
$12$, soit $4$ faces sur $12$, qui ont la même probabilité.}{équiprobabilité}
\begin{calculs}
p &= \dfrac{4}{12} &&\justif{$4$ faces favorables sur $12$}\\
p &= \dfrac{\barre{4}\times 1}{\barre{4}\times 3} &&\justif{on simplifie par $4$}\\
p &= \dfrac{1}{3} &&
\end{calculs}
\explique{Donc \res{$Y$ suit la loi de Bernoulli de paramètre $p=\dfrac{1}{3}$}.}{$Y$ vaut
$1$ (succès) ou $0$ (échec)}
\begin{center}
\renewcommand{\arraystretch}{2.4}
\begin{tabular}{|c|c|c|}
\hline
$y_i$ & $1$ & $0$\\
\hline
$P(Y=y_i)$ & $\dfrac{1}{3}$ & $\dfrac{2}{3}$\\
\hline
\end{tabular}
\end{center}
\explique{$E(Y)=p$, soit \res{$E(Y)=\dfrac{1}{3}$}.}{propriété 3}
\cas{c)}
\explique{$Z$ prend douze valeurs, de $1$ à $12$ : \res{$Z$ ne suit pas une loi de
Bernoulli}.}{une loi de Bernoulli n'a que les valeurs $1$ et $0$}
\end{exo}

\begin{exo}[Épreuve de Bernoulli ou non ?]
\cas{a)}
\explique{On note la taille : trois issues (S, M, L), donc \res{ce n'est pas
une épreuve de Bernoulli}.}{une épreuve de Bernoulli a exactement deux issues}
\cas{b)}
\explique{Deux issues : \res{c'est une épreuve de Bernoulli}, de succès \og la pièce tombe sur
Pile\fg{}, de probabilité \res{$\dfrac{1}{2}$}.}{pièce équilibrée}
\cas{c)}
\explique{On note le mois : douze issues, donc \res{ce n'est pas une épreuve de
Bernoulli}.}{plus de deux issues}
\cas{d)}
\explique{Deux issues (le colis arrive en bon état ou non) : \res{c'est une épreuve de
Bernoulli}, de succès \og le colis arrive en bon état\fg{} ; \res{sa probabilité n'est pas
connue}.}{l'énoncé ne donne pas la part des colis qui arrivent en bon état}
\cas{e)}
\explique{Deux issues (as ou non) : \res{c'est une épreuve de Bernoulli}, de succès \og la
carte est un as\fg{}.}{un jeu de $32$ cartes compte $4$ as}
\begin{calculs}
p &= \dfrac{4}{32} &&\justif{$4$ as sur $32$ cartes}\\
p &= \dfrac{\barre{4}\times 1}{\barre{4}\times 8} &&\justif{on simplifie par $4$}\\
p &= \res{\dfrac{1}{8}} &&
\end{calculs}
\rem{Dans a) et c), on obtiendrait une épreuve de Bernoulli en ne regardant qu'une
question à deux réponses (\og le vélo est-il de taille M ?\fg{}) : c'est le choix de ce
qu'on observe qui compte.}
\end{exo}

\begin{exo}[Le défi du dodécaèdre]
\cas{1)}
\explique{Les multiples de $4$ jusqu'à $12$ sont $4$ ; $8$ ; $12$, soit $3$ faces sur $12$, qui
ont la même probabilité.}{dé équilibré}
\begin{calculs}
P(S) &= \dfrac{3}{12} &&\justif{$3$ faces favorables sur $12$}\\
P(S) &= \dfrac{\barre{3}\times 1}{\barre{3}\times 4} &&\justif{on simplifie par $3$}\\
P(S) &= \res{\dfrac{1}{4}} &&\justif{soit $0{,}25$}
\end{calculs}
\cas{2)}
\explique{$X$ suit la loi de Bernoulli de paramètre $p=0{,}25$.}{deux issues : $S$ (succès)
ou non}
\begin{center}
\renewcommand{\arraystretch}{1.4}
\begin{tabular}{|c|c|c|}
\hline
$x_i$ & $1$ & $0$\\
\hline
$P(X=x_i)$ & $0{,}25$ & $0{,}75$\\
\hline
\end{tabular}
\end{center}
\explique{$E(X)=p$, soit \res{$E(X)=0{,}25$}.}{propriété 3}
\cas{3)}
\explique{Si $S$ est réalisé, $G=7-2=5$ ; sinon, $G=-2$.}{on déduit le prix de la partie}
\begin{center}
\renewcommand{\arraystretch}{1.4}
\begin{tabular}{|c|c|c|}
\hline
$g_i$ & $-2$ & $5$\\
\hline
$P(G=g_i)$ & $0{,}75$ & $0{,}25$\\
\hline
\end{tabular}
\end{center}
\explique{\res{$G$ ne suit pas une loi de Bernoulli} : elle a deux valeurs, mais ce sont $-2$
et $5$, pas $0$ et $1$.}{une loi de Bernoulli prend les valeurs $1$ et $0$}
\cas{4)}
\begin{calculs}
E(G) &= (-2)\times 0{,}75+5\times 0{,}25 &&\justif{valeur $\times$ probabilité}\\
E(G) &= -1{,}50+1{,}25 &&\\
E(G) &= \res{-0{,}25} &&
\end{calculs}
\explique{Sur un grand nombre de parties, le joueur \res{perd en moyenne $0{,}25$~€ par
partie}.}{$E(G)<0$ : le jeu est défavorable au joueur}
\end{exo}

\partie{Répétition d'épreuves de Bernoulli}

\begin{exo}[Calculer une probabilité avec un arbre]
\cas{a)}
\begin{minipage}[c]{0.52\linewidth}\raggedright
On note $R$ \og la commande est livrée à l'heure\fg{} : $P(R)=0{,}8$ et
$P(\overline{R})=1-0{,}8$, soit \res{$P(\overline{R})=0{,}2$}. Au bout de chaque chemin, la valeur de $X$.
\end{minipage}\hfill
\begin{minipage}[c]{0.44\linewidth}\centering
\begin{tikzpicture}
  \arbredeux{R}{0{,}8}{0{,}2}
  \feuille{SS}{$X=2$}
  \feuille{SE}{$X=1$}
  \feuille{ES}{$X=1$}
  \feuille{EE}{$X=0$}
\end{tikzpicture}
\end{minipage}
\begin{calculs}
P(X=2) &= 0{,}8\times 0{,}8 &&\justif{un seul chemin : $(R\,;R)$}\\
P(X=2) &= \res{0{,}64} &&
\end{calculs}
\begin{calculs}
P(X=1) &= 0{,}8\times 0{,}2+0{,}2\times 0{,}8 &&\justif{deux chemins :
$(R\,;\overline{R})$ et $(\overline{R}\,;R)$}\\
P(X=1) &= 0{,}16+0{,}16 &&\\
P(X=1) &= \res{0{,}32} &&
\end{calculs}
\cas{b)}
\begin{minipage}[c]{0.46\linewidth}\raggedright
On note $B$ \og le colis est remis du premier coup\fg{} : $P(B)=0{,}9$ et $P(\overline{B})=1-0{,}9$,
soit \res{$P(\overline{B})=0{,}1$}. Au bout de chaque chemin, la valeur de $Y$.
\end{minipage}\hfill
\begin{minipage}[c]{0.50\linewidth}\centering
\begin{tikzpicture}[y=0.75cm]
  \arbretrois{B}{0{,}9}{0{,}1}
  \feuille{SSS}{$Y=3$} \feuille{SSE}{$Y=2$} \feuille{SES}{$Y=2$}
  \feuille{SEE}{$Y=1$} \feuille{ESS}{$Y=2$} \feuille{ESE}{$Y=1$}
  \feuille{EES}{$Y=1$} \feuille{EEE}{$Y=0$}
\end{tikzpicture}
\end{minipage}
\begin{calculs}
P(Y=3) &= 0{,}9\times 0{,}9\times 0{,}9 &&\justif{un seul chemin : $(B\,;B\,;B)$}\\
P(Y=3) &= \res{0{,}729} &&
\end{calculs}
\begin{calculs}
P(Y=2) &= 0{,}081+0{,}081+0{,}081 &&\justif{trois chemins, chacun
$0{,}9\times 0{,}9\times 0{,}1$}\\
P(Y=2) &= \res{0{,}243} &&
\end{calculs}
\begin{calculs}
P(Y\ge 1) &= 1-P(Y=0) &&\justif{contraire : aucun colis remis du premier coup}\\
P(Y\ge 1) &= 1-0{,}1\times 0{,}1\times 0{,}1 &&\justif{le chemin
$(\overline{B}\,;\overline{B}\,;\overline{B})$}\\
P(Y\ge 1) &= 1-0{,}001 &&\\
P(Y\ge 1) &= \res{0{,}999} &&
\end{calculs}
\rem{Pour \og au moins un\fg{}, l'événement contraire n'utilise qu'un chemin, au lieu de
sept.}
\end{exo}

\begin{exo}[Trois lancers d'une pièce]
\cas{1)}%
\begin{minipage}[c]{0.46\linewidth}\raggedright
On note $P$ \og la pièce tombe sur Pile\fg{} ; $\overline{P}$ est \og Face\fg{}. La pièce
est équilibrée : chaque branche a la probabilité $0{,}5$. Au bout de chaque chemin, la
valeur de $X$.
\end{minipage}\hfill
\begin{minipage}[c]{0.50\linewidth}\centering
\begin{tikzpicture}[y=0.75cm]
  \arbretrois{P}{0{,}5}{0{,}5}
  \feuille{SSS}{$X=3$} \feuille{SSE}{$X=2$} \feuille{SES}{$X=2$}
  \feuille{SEE}{$X=1$} \feuille{ESS}{$X=2$} \feuille{ESE}{$X=1$}
  \feuille{EES}{$X=1$} \feuille{EEE}{$X=0$}
\end{tikzpicture}
\end{minipage}
\cas{2)}
\explique{Chaque chemin a la probabilité $0{,}5\times 0{,}5\times 0{,}5=0{,}125$.}{les $8$
chemins ont la même probabilité}
\begin{calculs}
P(X=3) &= \res{0{,}125} &&\justif{un seul chemin : $(P\,;P\,;P)$}
\end{calculs}
\begin{calculs}
P(X=2) &= 3\times 0{,}125 &&\justif{trois chemins avec deux Pile}\\
P(X=2) &= \res{0{,}375} &&
\end{calculs}
\cas{3)}
\explique{\og Au moins une fois Face\fg{}, c'est \og pas trois Pile\fg{} : \res{$\{X\le
2\}$}.}{contraire de $\{X=3\}$}
\begin{calculs}
P(X\le 2) &= 1-P(X=3) &&\justif{événement contraire}\\
P(X\le 2) &= 1-0{,}125 &&\\
P(X\le 2) &= \res{0{,}875} &&
\end{calculs}
\cas{4)}
\explique{$P(X=0)=0{,}125$ (un chemin) et $P(X=1)=3\times 0{,}125=0{,}375$ (trois
chemins).}{comme à la question 2}
\begin{center}
\renewcommand{\arraystretch}{1.4}
\begin{tabular}{|c|c|c|c|c|}
\hline
$x_i$ & $0$ & $1$ & $2$ & $3$\\
\hline
$P(X=x_i)$ & $0{,}125$ & $0{,}375$ & $0{,}375$ & $0{,}125$\\
\hline
\end{tabular}
\end{center}
\explique{Vérification : $0{,}125+0{,}375+0{,}375+0{,}125=1$.}{propriété 1}
\begin{calculs}
E(X) &= 0\times 0{,}125+1\times 0{,}375+2\times 0{,}375+3\times 0{,}125 &&\justif{valeur
$\times$ probabilité}\\
E(X) &= 0+0{,}375+0{,}75+0{,}375 &&\\
E(X) &= \res{1{,}5} &&
\end{calculs}
\explique{Sur un grand nombre de séries de trois lancers, on obtient en moyenne $1{,}5$
Pile.}{la moitié de $3$ lancers : la pièce est équilibrée}
\end{exo}

\begin{exo}[Une agence de location]
\cas{1)}%
\begin{minipage}[c]{0.46\linewidth}\raggedright
On note $A$ \og la voiture est en entretien\fg{} : $P(A)=0{,}2$ et
$P(\overline{A})=0{,}8$. Les voitures sont indépendantes : les mêmes probabilités à chaque niveau. Au bout de
chaque chemin, la valeur de $X$.
\end{minipage}\hfill
\begin{minipage}[c]{0.50\linewidth}\centering
\begin{tikzpicture}[y=0.75cm]
  \arbretrois{A}{0{,}2}{0{,}8}
  \feuille{SSS}{$X=3$} \feuille{SSE}{$X=2$} \feuille{SES}{$X=2$}
  \feuille{SEE}{$X=1$} \feuille{ESS}{$X=2$} \feuille{ESE}{$X=1$}
  \feuille{EES}{$X=1$} \feuille{EEE}{$X=0$}
\end{tikzpicture}
\end{minipage}
\cas{2)}
\begin{calculs}
P(X=3) &= 0{,}2\times 0{,}2\times 0{,}2 = \res{0{,}008} &&\justif{un chemin :
$(A\,;A\,;A)$}\\
P(X=2) &= 3\times 0{,}2\times 0{,}2\times 0{,}8 = \res{0{,}096} &&\justif{trois chemins,
chacun $0{,}032$}\\
P(X=1) &= 3\times 0{,}2\times 0{,}8\times 0{,}8 = \res{0{,}384} &&\justif{trois chemins,
chacun $0{,}128$}\\
P(X=0) &= 0{,}8\times 0{,}8\times 0{,}8 = \res{0{,}512} &&\justif{un chemin :
$(\overline{A}\,;\overline{A}\,;\overline{A})$}
\end{calculs}
\begin{center}
\renewcommand{\arraystretch}{1.4}
\begin{tabular}{|c|c|c|c|c|}
\hline
$x_i$ & $0$ & $1$ & $2$ & $3$\\
\hline
$P(X=x_i)$ & $0{,}512$ & $0{,}384$ & $0{,}096$ & $0{,}008$\\
\hline
\end{tabular}
\end{center}
\explique{Vérification : $0{,}512+0{,}384+0{,}096+0{,}008=1$.}{propriété 1}
\cas{3)}
\begin{calculs}
E(X) &= 0\times 0{,}512+1\times 0{,}384+2\times 0{,}096+3\times 0{,}008 &&\justif{valeur
$\times$ probabilité}\\
E(X) &= 0+0{,}384+0{,}192+0{,}024 &&\\
E(X) &= \res{0{,}6} &&
\end{calculs}
\explique{$3\times 0{,}2=0{,}6$ : \res{on trouve la même valeur}. Sur $3$ voitures, il y en
a en moyenne $0{,}6$ en entretien.}{chaque voiture compte en moyenne pour $0{,}2$ :
l'espérance d'une loi de Bernoulli de paramètre $0{,}2$}
\cas{4)}
\begin{calculs}
P(X\le 1) &= P(X=0)+P(X=1) &&\justif{au plus une voiture : $0$ ou $1$}\\
P(X\le 1) &= 0{,}512+0{,}384 &&\\
P(X\le 1) &= \res{0{,}896} &&
\end{calculs}
\end{exo}

\begin{exo}[Fléchettes électroniques]
\cas{1)}%
\begin{minipage}[c]{0.52\linewidth}\raggedright
On note $R$ \og la fléchette atteint la zone rouge\fg{} : $P(R)=0{,}6$ et
$P(\overline{R})=0{,}4$. Au bout de chaque chemin, le score : $2+2=4$, $2-1=1$, $-1+2=1$,
$-1-1=-2$.
\end{minipage}\hfill
\begin{minipage}[c]{0.44\linewidth}\centering
\begin{tikzpicture}
  \arbredeux{R}{0{,}6}{0{,}4}
  \feuille{SS}{$X=4$}
  \feuille{SE}{$X=1$}
  \feuille{ES}{$X=1$}
  \feuille{EE}{$X=-2$}
\end{tikzpicture}
\end{minipage}
\cas{2)}
\explique{$X$ prend les valeurs \res{$-2$ ; $1$ et $4$}.}{deux échecs, un succès, deux succès}
\cas{3)}
\begin{calculs}
P(X=4) &= 0{,}6\times 0{,}6 = \res{0{,}36} &&\justif{chemin $(R\,;R)$}\\
P(X=1) &= 2\times 0{,}6\times 0{,}4 = \res{0{,}48} &&\justif{chemins $(R\,;\overline{R})$ et
$(\overline{R}\,;R)$}\\
P(X=-2) &= 0{,}4\times 0{,}4 = \res{0{,}16} &&\justif{chemin
$(\overline{R}\,;\overline{R})$}
\end{calculs}
\begin{center}
\renewcommand{\arraystretch}{1.4}
\begin{tabular}{|c|c|c|c|}
\hline
$x_i$ & $-2$ & $1$ & $4$\\
\hline
$P(X=x_i)$ & $0{,}16$ & $0{,}48$ & $0{,}36$\\
\hline
\end{tabular}
\end{center}
\explique{Vérification : $0{,}16+0{,}48+0{,}36=1$.}{propriété 1}
\cas{4)}
\begin{calculs}
P(X>0) &= P(X=1)+P(X=4) &&\justif{les scores positifs}\\
P(X>0) &= 0{,}48+0{,}36 &&\\
P(X>0) &= \res{0{,}84} &&
\end{calculs}
\explique{Le joueur \res{termine avec un score positif dans $84\,\%$ des cas}, sur un grand
nombre de séries de deux lancers.}{il suffit qu'une fléchette sur deux atteigne la zone
rouge}
\cas{5)}
\begin{calculs}
E(X) &= (-2)\times 0{,}16+1\times 0{,}48+4\times 0{,}36 &&\justif{valeur $\times$
probabilité}\\
E(X) &= -0{,}32+0{,}48+1{,}44 &&\\
E(X) &= \res{1{,}6} &&
\end{calculs}
\explique{Sur un grand nombre de séries de deux lancers, le joueur \res{marque en moyenne
$1{,}6$ point}.}{propriété 2}
\end{exo}

\partie{Simulation d'échantillons}

\begin{exo}[Simuler à la main et à l'ordinateur]
\cas{1)}
\explique{\texttt{if random() < \res{0.8}:} puis \texttt{return \res{0}}.}{\texttt{random()}
est inférieur à $0{,}8$ avec la probabilité $0{,}8$ : jour à plus de $50$ ventes ; sinon,
échec}
\cas{2)}
\begin{calculs}
f_{1} &= \dfrac{8}{10} &&\justif{chiffres $8$ et $9$, donc $2$ échecs et $8$ jours sur $10$}\\
f_{1} &= \res{0{,}8} &&\\
f_{2} &= \dfrac{7}{10} &&\justif{chiffres $9$, $8$ et $9$, donc $3$ échecs et $7$ jours}\\
f_{2} &= \res{0{,}7} &&
\end{calculs}
\cas{3)}
\begin{calculs}
f &= \dfrac{16}{20} &&\justif{$16$ jours sur $20$}\\
f &= \res{0{,}8} &&\justif{$\frac{16\times 5}{20\times 5}=\frac{80}{100}$}
\end{calculs}
\explique{Chaque échantillon est tiré au hasard : \res{la fréquence change d'un échantillon
à l'autre}, autour de $0{,}8$.}{fluctuation d'échantillonnage (partie VII)}
\cas{4)}
\begin{calculs}
\text{total} &= 16+15+17+16+14+18+16+15+17+16 &&\justif{jours dans les $10$ échantillons}\\
\text{total} &= 160 &&\justif{sur $10\times 20=200$ jours}\\
\text{moyenne des fréquences} &= \dfrac{160}{200} &&\\
\text{moyenne des fréquences} &= \res{0{,}8} &&
\end{calculs}
\explique{La moyenne des fréquences est \res{égale à $p=0{,}8$} ici.}{les fréquences
fluctuent autour de $p$}
\end{exo}

\begin{exo}[Lire un programme]
\cas{1)}
\explique{$5\,\%$ des livres ont un défaut : \res{$P(D)=0{,}05$}.}{un pourcentage écrit en
décimal}
\cas{2)}
\explique{\texttt{livre()} \res{renvoie $1$ si le livre simulé a un défaut, $0$ sinon} ; elle
renvoie $1$ quand \texttt{random() < 0.05}, donc \res{avec la probabilité $0{,}05$}.}{elle
simule la loi de Bernoulli de paramètre $0{,}05$}
\cas{3)}
\begin{calculs}
f &= \dfrac{493}{10\,000} &&\justif{nombre de livres avec un défaut sur $10\,000$}\\
f &= \res{0{,}0493} &&
\end{calculs}
\cas{4)}
\explique{\res{Ce n'est pas étonnant} : sur un grand échantillon ($10\,000$ livres), la
fréquence observée $0{,}0493$ est très proche de $p=0{,}05$.}{plus l'échantillon est
grand, moins la fréquence fluctue (partie VII)}
\end{exo}

\begin{exo}[La moyenne des fréquences]
\cas{1)}
\begin{calculs}
\text{somme} &= 0{,}28+0{,}33+0{,}31+0{,}27+0{,}31 &&\justif{on additionne les $5$ fréquences}\\
\text{somme} &= 1{,}5 &&\\
\text{moyenne} &= \dfrac{1{,}5}{5} &&\justif{on divise par leur nombre}\\
\text{moyenne} &= \res{0{,}3} &&
\end{calculs}
\cas{2)}
\explique{Deux issues ; le succès est \og le client choisit l'option\fg{}, de probabilité
$0{,}3$ : \res{$X$ suit la loi de Bernoulli de paramètre $0{,}3$}, et \res{$E(X)=0{,}3$}. La
moyenne des fréquences \res{est égale à l'espérance} ici.}{propriété 3 : $E(X)=p$}
\cas{3)}
\explique{La moyenne des fréquences de $500$ échantillons vaut $0{,}2997$ : \res{les
fréquences fluctuent autour de $0{,}3$}, et leur moyenne en est très proche.}{certaines
fréquences sont au-dessus de $0{,}3$, d'autres en dessous}
\end{exo}

\partie{Fluctuation d'échantillonnage}

\begin{exo}[Fréquente ou rare ?]
\cas{1)}
\begin{calculs}
p-2s &= 0{,}5-2\times 0{,}05 = \res{0{,}40} &&\justif{$2s=0{,}10$}\\
p+2s &= 0{,}5+2\times 0{,}05 = \res{0{,}60} &&
\end{calculs}
\cas{2)}
\explique{$f=\dfrac{57}{100}=0{,}57$, qui est dans $[0{,}40\,;0{,}60]$ : c'est une
\res{situation fréquente}. Rien d'étonnant pour une pièce équilibrée.}{distance à $p$ :
$0{,}07$, inférieure à $2s=0{,}10$}
\cas{3)}
\explique{$f=\dfrac{68}{100}=0{,}68$, qui est hors de $[0{,}40\,;0{,}60]$ : c'est une
\res{situation rare} dans le modèle. On peut penser que \res{cette autre pièce
n'est pas équilibrée} : elle tombe plus souvent sur Pile.}{distance à $p$ : $0{,}18$,
supérieure à $2s=0{,}10$}
\cas{4)}
\explique{$400=4\times 100$ : la taille est multipliée par $4$, donc $s$ est divisé par $2$ :
\res{$s\approx 0{,}025$}.}{propriété 5}
\end{exo}

\begin{exo}[Un générateur de chiffres]
\cas{1)}
\explique{On lit la barre au-dessus de $12$ : \res{$29$ paquets}.}{hauteur de la barre}
\begin{calculs}
f &= \dfrac{12}{25} &&\justif{$12$ chiffres égaux à $1$ sur $25$}\\
f &= \res{0{,}48} &&\justif{$\frac{12\times 4}{25\times 4}=\frac{48}{100}$}
\end{calculs}
\cas{2)}
\begin{calculs}
p-2s &= 0{,}5-2\times 0{,}1 = \res{0{,}3} &&\justif{$2s=0{,}2$}\\
p+2s &= 0{,}5+2\times 0{,}1 = \res{0{,}7} &&
\end{calculs}
\explique{En nombre de $1$ sur $25$, on a $0{,}3\times 25=7{,}5$ et $0{,}7\times 25=17{,}5$. La
fréquence est dans $[0{,}3\,;0{,}7]$ pour \res{$8$ à $17$ chiffres égaux à $1$}.}{$7$ donne
$0{,}28$, trop petit ; $18$ donne $0{,}72$, trop grand}
\cas{3)}
\begin{calculs}
\text{paquets hors de l'intervalle} &= 1+1+6+2 &&\justif{barres de $6$ ; $7$ ; $18$ ;
$19$}\\
\text{paquets hors de l'intervalle} &= \res{10} &&
\end{calculs}
\begin{calculs}
\text{pourcentage} &= \dfrac{10}{200}\times 100 &&\justif{part des $200$ paquets}\\
\text{pourcentage} &= \res{5} &&
\end{calculs}
\explique{$5\,\%$ des paquets sont hors de $[p-2s\,;p+2s]$, donc $95\,\%$ dedans :
\res{c'est conforme au cours}.}{propriété 5 : environ $95\,\%$ des fréquences dans
$[p-2s\,;p+2s]$}
\cas{4)}
\begin{calculs}
f &= \dfrac{21}{25} &&\justif{$21$ chiffres égaux à $1$ sur $25$}\\
f &= 0{,}84 &&\justif{$\frac{21\times 4}{25\times 4}=\frac{84}{100}$}
\end{calculs}
\explique{$0{,}84$ est hors de $[0{,}3\,;0{,}7]$ : c'est une \res{situation rare} dans le
modèle ; on peut douter du générateur.}{distance à $p$ : $0{,}34$, supérieure à $2s=0{,}2$ ;
aucun des $200$ paquets simulés n'en avait autant}
\end{exo}

\begin{exo}[La taille de l'échantillon]
\cas{1)}
\explique{De $n=25$ à $n=100$, $n$ est multiplié par $4$ et $s$ passe de $0{,}10$ à $0{,}05$ :
\res{$s$ est divisé par $2$}. Pour $n=400=4\times 100$ : \res{$s\approx 0{,}025$}.}{propriété
5}
\cas{2)}
\begin{calculs}
\dfrac{2\,500}{25} &= \res{100} &&\justif{la taille est multipliée par $100$}
\end{calculs}
\explique{Quand $n$ est multiplié par $100$, $s$ est divisé par $10$, donc $\frac{0{,}10}{10}$,
soit \res{$s\approx 0{,}01$}.}{propriété 5}
\cas{3)}
\explique{Avec $2\,500$ usagers, $s\approx 0{,}01$ au lieu de $0{,}05$ : \res{la fréquence
observée fluctue beaucoup moins, elle est plus proche de la vraie part des usagers
satisfaits}.}{l'écart-type est divisé par $5$}
\end{exo}

\begin{exo}[Des batteries longue durée]
\cas{1)}
\explique{Deux issues ; le succès est \og la batterie supporte plus de $1\,000$
recharges\fg{}, de probabilité $0{,}9$ : \res{$X$ suit la loi de Bernoulli de paramètre
$0{,}9$}, et \res{$E(X)=0{,}9$}.}{propriété 3 : $E(X)=p$}
\cas{2)}
\begin{calculs}
p-2s &= 0{,}9-2\times 0{,}03 = \res{0{,}84} &&\justif{$2s=0{,}06$}\\
p+2s &= 0{,}9+2\times 0{,}03 = \res{0{,}96} &&
\end{calculs}
\cas{3)}
\explique{$f=\dfrac{87}{100}=0{,}87$, qui est dans $[0{,}84\,;0{,}96]$ : c'est une
\res{situation fréquente}, conforme à l'annonce.}{distance à $p$ : $0{,}03$, inférieure à
$2s=0{,}06$}
\cas{4)}
\explique{$f=\dfrac{78}{100}=0{,}78$, qui est hors de $[0{,}84\,;0{,}96]$ : c'est une
\res{situation rare} dans le modèle. On peut penser que \res{ce lot compte moins de
$90\,\%$ de batteries longue durée}.}{distance à $p$ : $0{,}12$, supérieure à $2s=0{,}06$}
\cas{5)}
\explique{\res{Non} : une situation rare reste possible par hasard ; le laboratoire a
de bonnes raisons de douter du lot, mais pas de certitude.}{le hasard peut produire un
échantillon éloigné de $p$, même si c'est peu probable}
\end{exo}

\besoinplace{16\baselineskip}
\partiesansnum{Exercices type épreuve anticipée}

\begin{exo}[QCM]
\cas{1.}
\begin{calculs}
P(X=2) &= 1-(0{,}2+0{,}5) &&\justif{la somme des probabilités vaut $1$}\\
P(X=2) &= 1-0{,}7 = 0{,}3 &&
\end{calculs}
\explique{Réponse \res{b)}.}{a) est $P(X=0)+P(X=1)$, pas $P(X=2)$}
\cas{2.}
\begin{calculs}
P(X\ge 5) &= P(X=5)+P(X=6) &&\justif{les faces $5$ et $6$}\\
P(X\ge 5) &= 0{,}2+0{,}35 = 0{,}55 &&
\end{calculs}
\explique{Réponse \res{c)}.}{a) suppose le dé équilibré : il est truqué}
\cas{3.}
\begin{calculs}
E(X) &= 0\times 0{,}5+2\times 0{,}3+10\times 0{,}2 &&\justif{valeur $\times$ probabilité}\\
E(X) &= 0+0{,}6+2 = 2{,}6 &&
\end{calculs}
\explique{Réponse \res{a)}.}{b) est la moyenne des valeurs $0$, $2$ et $10$, sans tenir
compte des probabilités}
\cas{4.}
\explique{$E(X)=p=0{,}3$. Réponse \res{c)}.}{propriété 3}
\cas{5.}
\begin{calculs}
P &= 0{,}8\times 0{,}8 = 0{,}64 &&\justif{un seul chemin de l'arbre : carte, puis carte}
\end{calculs}
\explique{Réponse \res{b)}.}{a) additionne au lieu de multiplier : plus grand que $1$}
\cas{6.}
\explique{$p-2s=0{,}40$ et $p+2s=0{,}60$ ; $0{,}62$ est hors de $[0{,}40\,;0{,}60]$ : une
situation rare. Réponse \res{b)}.}{$2s=0{,}10$ ; d) est faux : rare ne veut pas dire
impossible}
\end{exo}

\begin{exo}[Vrai ou faux ?]
\cas{1.}
\begin{calculs}
P(X\ge 2) &= 0{,}3+0{,}2+0{,}1 &&\justif{les valeurs $2$ ; $3$ ; $4$}\\
P(X\ge 2) &= 0{,}6 &&
\end{calculs}
\explique{$0{,}6\neq 0{,}3$ : l'affirmation 1 est \res{fausse}.}{$0{,}3$ est $P(X=2)$
seulement}
\cas{2.}
\begin{calculs}
E(X) &= 1\times 0{,}4+2\times 0{,}3+3\times 0{,}2+4\times 0{,}1 &&\justif{valeur $\times$
probabilité}\\
E(X) &= 0{,}4+0{,}6+0{,}6+0{,}4 &&\\
E(X) &= 2 &&
\end{calculs}
\explique{En moyenne, un client achète $2$ articles : l'affirmation 2 est
\res{vraie}.}{propriété 2}
\cas{3.}
\begin{calculs}
P(X\le 2) &= 0{,}4+0{,}3 &&\justif{les valeurs $1$ et $2$}\\
P(X\le 2) &= 0{,}7 &&
\end{calculs}
\explique{$70\,\%$ des clients achètent au plus $2$ articles, et $70\,\%>50\,\%$ :
l'affirmation 3 est \res{vraie}.}{}
\cas{4.}
\begin{calculs}
P &= 0{,}4\times 0{,}4 = 0{,}16 &&\justif{deux choix indépendants : on multiplie}
\end{calculs}
\explique{L'affirmation 4 est \res{vraie}.}{un seul chemin de l'arbre : $X=1$, puis $X=1$}
\end{exo}

\begin{exo}[Les nuits d'un hôtel]
\cas{1.}
\explique{$X$ prend les valeurs $1$ à $5$ ; chaque pourcentage s'écrit en
décimal.}{$30\,\%=0{,}30$}
\begin{center}
\renewcommand{\arraystretch}{1.4}
\begin{tabular}{|c|c|c|c|c|c|}
\hline
$x_i$ & $1$ & $2$ & $3$ & $4$ & $5$\\
\hline
$P(X=x_i)$ & $0{,}30$ & $0{,}35$ & $0{,}20$ & $0{,}10$ & $0{,}05$\\
\hline
\end{tabular}
\end{center}
\explique{Vérification : $0{,}30+0{,}35+0{,}20+0{,}10+0{,}05=\res{1}$.}{propriété 1}
\cas{2.}
\begin{calculs}
P(X\ge 2) &= 1-P(X=1) &&\justif{contraire : le client ne reste qu'une nuit}\\
P(X\ge 2) &= 1-0{,}30 &&\\
P(X\ge 2) &= \res{0{,}70} &&
\end{calculs}
\explique{\res{$70\,\%$ des clients restent au moins $2$ nuits} à l'hôtel.}{on vérifie :
$0{,}35+0{,}20+0{,}10+0{,}05=0{,}70$}
\cas{3.}
\begin{calculs}
E(X) &= 1\times 0{,}30+2\times 0{,}35+3\times 0{,}20+4\times 0{,}10+5\times 0{,}05
  &&\justif{valeur $\times$ probabilité}\\
E(X) &= 0{,}30+0{,}70+0{,}60+0{,}40+0{,}25 &&\\
E(X) &= \res{2{,}25} &&
\end{calculs}
\explique{Sur un grand nombre de clients, un client \res{passe en moyenne $2{,}25$ nuits} à
l'hôtel.}{propriété 2}
\cas{4.}
\begin{calculs}
\text{recette moyenne} &= 80\times 2{,}25 &&\justif{prix d'une nuit $\times$ nombre moyen de
nuits}\\
\text{recette moyenne} &= \res{180} &&\justif{$80\times 2=160$ et $80\times 0{,}25=20$}
\end{calculs}
\explique{L'hôtel peut attendre une recette moyenne de \res{$180$~€ par client}.}{}
\cas{5.}
\explique{On note $U$ \og le client ne reste qu'une nuit\fg{} : $P(U)=0{,}3$. Les deux
choix sont indépendants : sur l'arbre, on suit le chemin $(U\,;U)$.}{arbre à deux niveaux,
branches $U$ et $\overline{U}$}
\begin{calculs}
P &= 0{,}3\times 0{,}3 &&\justif{produit des probabilités du chemin}\\
P &= \res{0{,}09} &&
\end{calculs}
\explique{Les deux clients ne restent qu'une nuit avec la probabilité $0{,}09$, soit
$9\,\%$.}{}
\end{exo}

\end{document}
